\chapter{Structural VAR}
\label{cap:svar}
% ── index ───────────────────────────────────────────────────────
\index{structural VAR|textbf}
\index{SVAR|see{structural VAR}}
\index{svar@\texttt{svar()}|textbf}
\index{sirf@\texttt{sirf()}|textbf}
\index{Cholesky, decomposition of|textbf}
\index{structural identification|textbf}
\index{structural impulse-response function|textbf}
\index{causal ordering}

% ─────────────────────────────────────────────────────────────────────────────
\section{From reduced form to structure}

The reduced-form VAR (ch.~\ref{cap:var}) provides $K^2 \cdot p + K$ free parameters.
The reduced-form impulse response measures the effect of \emph{innovations} $u_t$
with no economic content. The SVAR recovers \emph{structural shocks}
$\varepsilon_t$ (supply, demand, monetary policy) by imposing restrictions
that reflect economic theory:

\[
  A\,y_t = A_1^*\,y_{t-1} + \cdots + B\,\varepsilon_t
  \qquad
  \varepsilon_t \sim (0, I_K)
\]

\subsection*{Cholesky Identification}

The lower triangular decomposition of $\Sigma_u = P\,P'$ imposes a
recursive causal ordering: the variable in position 1 is contemporaneously
exogenous; the one in position 2 responds contemporaneously only
to the first; etc.

\[
  B = P \quad\text{(lower triangular)} \implies y_{1,t} \perp \varepsilon_{2,t} \text{ in the same period}
\]

The structural IRF $\Theta_h = \Phi_h B$ measures the response of $y_{k,t+h}$
to a unit shock in $\varepsilon_{j,t}$.

% ─────────────────────────────────────────────────────────────────────────────
\section{DGP Verification}

DGP: bivariate VAR(1) with true Cholesky restriction (y2 responds
contemporaneously to y1, but not vice versa).

\begin{lstlisting}[caption={Structural DGP and Cholesky SVAR}]
seed(42)
generate df t  = _n
// VAR(1): y1 and y2 with cross coefficients
...
replace df y1 = 0.5*ly1 + 0.3*ly2 + e1t          if t == it
replace df y2 = 0.2*ly1 + 0.4*ly2 + 0.6*e1t + e2t if t == it

let m_var  = var(df, y1, y2, lags=1)
display m_var

let m_svar = svar(df, y1, y2, lags=1, id=cholesky)
display m_svar

let m_sirf = sirf(m_svar, steps=10)
display m_sirf
\end{lstlisting}

\subsection*{Reduced-form VAR}

\begin{lstlisting}[language={}, basicstyle=\ttfamily\small,
  frame=none, numbers=none, backgroundcolor=\color{codbg},
  xleftmargin=0pt, xrightmargin=0pt, breaklines=false]
Vector Autoregression (VAR(1)):  Observations=299  AIC=-2.83  BIC=-2.76
Residual Covariance (Σ_u):
[ 0.2452   0.1617 ]
[ 0.1617   0.3375 ]
\end{lstlisting}

\subsection*{SVAR — B Matrix}

\begin{lstlisting}[language={}, basicstyle=\ttfamily\small,
  frame=none, numbers=none, backgroundcolor=\color{codbg},
  xleftmargin=0pt, xrightmargin=0pt, breaklines=false]
============================= SVAR(1) — Cholesky =============================
B Matrix (structural impact):
[ 0.4952   0.0000 ]     <- shock 1 affects y1, does not affect y2 contemporaneously
[ 0.3265   0.4805 ]     <- shock 2 affects y2; y2 responds to shock 1
\end{lstlisting}

The zero column in $B_{12}$ reflects the Cholesky ordering: shock 2 does not affect
$y_1$ contemporaneously. $B_{21} = 0{,}33$ estimates the contemporaneous effect of
shock 1 on $y_2$ (true DGP: 0{,}6 $\times$ $\sigma_{\varepsilon_1}$).

\subsection*{Structural IRF}

\begin{lstlisting}[language={}, basicstyle=\ttfamily\small,
  frame=none, numbers=none, backgroundcolor=\color{codbg},
  xleftmargin=0pt, xrightmargin=0pt, breaklines=false]
Impulse: Shock 1 (Var0)
  h=0:  y1=+0.495  y2=+0.327   (contemporaneous impact)
  h=1:  y1=+0.324  y2=+0.224
  h=5:  y1=+0.063  y2=+0.044   (progressive dampening)
  h=9:  y1=+0.012  y2=+0.009

Impulse: Shock 2 (Var1)
  h=0:  y1= 0.000  y2=+0.481   (Cholesky restriction: y1 does not respond)
  h=1:  y1=+0.115  y2=+0.170
\end{lstlisting}

The response of $y_1$ to shock 2 at $h=0$ is exactly zero — the Cholesky
identification restriction applied. At horizon $h=1$, $y_1$ begins to respond
via the lagged VAR channel.

% ─────────────────────────────────────────────────────────────────────────────
\section{Syntax}

\begin{lstlisting}
svar(df, y1, y2, lags=1, id=cholesky)     // SVAR with Cholesky identification
svar(df, y1, y2, lags=1, id=longrun)      // BQ: long-run restrictions
sirf(m_svar, steps=10)                    // structural IRF
\end{lstlisting}

\begin{center}
\begin{tabular}{ll}
\toprule
\textbf{Identification} & \textbf{Restrictions} \\
\midrule
\texttt{cholesky} & lower triangular — recursive short-run restrictions \\
\texttt{longrun}  & Blanchard-Quah — long-run restrictions \\
\bottomrule
\end{tabular}
\end{center}

\begin{nota}
  The ordering of variables in the Cholesky SVAR is crucial: the variable listed
  first is treated as contemporaneously exogenous to all others.
  The choice of ordering should be guided by economic theory or by
  results that are robust to permutations of the order.
\end{nota}
